% Pregunta 257 (2018-enero-OEP2017, n.º 13): Cournot con costes distintos. Autor: Víctor Gutiérrez Marcos
% A escala con A = 11, c_1 = 1, c_2 = 3:
%   R_1: q_1 = (10 - q_2)/2, de (5,0) a (0,10);  R_2: q_2 = (8 - q_1)/2, de (0,4) a (8,0)
%   Equilibrio E = (4, 2): q_1* - q_2* = c_2 - c_1 = 2; beneficios 16 y 4
\documentclass[tikz,border=6pt]{standalone}
\input{tikz-tcee}
\begin{document}
\begin{tikzpicture}[tcee,x=0.5cm,y=0.5cm]
  \ejes{$q_1$}{$q_2$}{11}{11}
  % Recta de 45 grados: cantidades iguales
  \draw[guia] (0,0) -- (9.5,9.5) node[etiqueta,above,text=tceegris] {$q_1=q_2$};
  % Función de reacción de E1 (coste bajo)
  \draw[curva1] (5,0) -- (0,10);
  \node[etiqueta,right,text=tceemorado] at (0.6,9.3) {$R_1(q_2)$};
  % Función de reacción de E2 (coste alto)
  \draw[curva2] (0,4) -- (8,0);
  \node[etiqueta,above,text=tceeazul] at (7.2,0.6) {$R_2(q_1)$};
  % Cortes con los ejes
  \node[etiqueta,left] at (0,10) {$10$};
  \node[etiqueta,left] at (0,4) {$4$};
  \node[etiqueta,below] at (5,0) {$5$};
  \node[etiqueta,below] at (8,0) {$8$};
  % Equilibrio de Cournot
  \node[punto] (E) at (4,2) {};
  \proyeccion{(E)}{$4$}{$2$}
  \node[etiqueta,above right] at (E) {$E$};
  \node[etiqueta,align=left,anchor=west] at (5.6,4.2) {$q_1^*=4$, \ $q_2^*=2$\\$\pi_1^*=16$, \ $\pi_2^*=4$};
\end{tikzpicture}
\end{document}
